\begin{textsample}{POS Dim 1 – ma – Score 46.00 – ma/ma\_ef\_51.txt}  \label{ex:f1_pos_065}
O componente curricular Geografia no Ensino Fundamental Para o ensino de Geografia, a BNCC propõe as mesmas unidades temáticas ao longo de todo o Ensino Fundamental, no entanto os objetos do conhecimento apresentam características diferenciadas entre os anos dessa etapa de ensino.
Isso se deve, entre outros aspectos, ao estágio de desenvolvimento cognitivo e socioespacial dos alunos, às \textbf{abordagens} dadas ao conhecimento geográfico nos diferentes anos, aos diferentes níveis de complexidade na \textbf{abordagem} \textbf{conceitual} e às diferentes escalas utilizadas – local, regional e mundial.
Geografia nos anos iniciais O objetivo das Ciências Humanas nos anos iniciais como componente curricular não é tão diferente dos demais níveis de ensino, embora tenha características \textbf{metodológicas} de acordo com a faixa etária dos alunos ( STRAFORINI, 2008 ).
Para \textbf{pesquisadores} dessa área, o ensino de Geografia para crianças é \textbf{um} dos espaços para a formação de cidadãos no cenário da realidade dela a partir da descoberta das desigualdades sociais desde a infância, de forma crítica.
Podemos afirmar \textbf{que} nos anos iniciais o aluno deve se reconhecer como \textbf{sujeito} para \textbf{que} ele construa sua visão de cidadão e se situe nos diferentes espaços e tempo \textbf{que} o cercam, experimentando o enfrentamento do seu cotidiano em constante transformação.
A Geografia é \textbf{um} instrumento dessa construção de conhecimento.
Nessa \textbf{direção} é \textbf{uma} alfabetização \textbf{que} está para além do “ domínio de técnicas de ler e escrever, mas com a amplitude de decifrar/decodificar outros signos diferentes de alfabetos ” ( MARTINS, 2015:66 ).
Nesse sentido, alfabetizar para as Ciências Humanas significa conhecimento da realidade local e global, construção da identidade e da noção de pertencimento ao fazer a leitura do mundo, das paisagens antigas e novas, dos instrumentos de comunicação, das novas relações de linguagem escrita, digital e informacional, ocorridas nos mais variados contextos sociais mudando forma e conteúdo ( PEREIRA, 2016 ).
Portanto, é resgatar a memória da sociedade, compreendendo como os \textbf{homens} se constituíram \textbf{historicamente} e como se relacionam socialmente entre si e com o ambiente, como resultado das diversas combinações de elementos do passado e do presente, constituindo novas espacialidades.
As relações com o contexto histórico e social são fundamentais para \textbf{que} o ensino seja eficaz, atraente, curioso e capaz de instrumentalizar os educandos à compreensão do mundo \textbf{atual}.
É \textbf{preciso} romper com \textbf{uma} visão fragmentada, linear das ações educativas e promover práticas \textbf{que} favoreçam a formação da cidadania ( MARTINS, 2015:65 ).
Porto ( 2015 ) \textbf{diz} \textbf{que} as primeiras leituras geográficas do mundo se iniciam nas vivências do cotidiano da criança, no experienciar diário de sua realidade física e simbólica no lugar onde \textbf{vive}, na geografia viva \textbf{que} está no seu cotidiano desde o nascimento, na cultura estabelecida \textbf{que} precede sua entrada na escola.
O ensino de Geografia nos anos iniciais do Ensino Fundamental deve ser \textbf{um} instrumento de “ olhar da criança para a escola, para o mundo, para a família, para a sala de aula em diferentes escalas ” ( PEREIRA, 2016 ).
Para tal, os Parâmetros Curriculares Nacionais \textbf{dizem} \textbf{que} “ a Geografia tem por objetivo estudar as relações entre o processo histórico na formação das sociedades humanas e o funcionamento da natureza por meio da leitura do lugar, do território, a partir de sua paisagem em diferentes noções espaciais e temporais ” ( BRASIL, 1998:26 ).
Callai ( 2005:229 ) \textbf{diz} \textbf{que} o lugar da Geografia nos anos iniciais é aprender a pensar a ler o espaço, e para Castellar ( apud CALLAI, 2005 ) é resultado de estimulação prática para compreender como os \textbf{homens} e mulheres se constituíram e se relacionam social e economicamente no espaço \textbf{vivido}.
Para dar conta dos objetivos da Geografia dos anos iniciais é necessário \textbf{que} a criança se aproprie dos conceitos da educação geográfica ( MARTINS, 2015 ).
Almeida ( 1999:830 ) \textbf{diz} \textbf{que} a finalidade da Geografia é “ munir os alunos de conhecimentos \textbf{que} lhes permitam agir de modo mais lúcido ao tratar das questões do espaço em diferentes níveis.
O ensino da Geografia tem, portanto, papel decisivo na formação da cidadania ”.
A questão é, portanto, qual é a importância da Geografia nos anos iniciais e como esse processo de aprendizagem contribui para o desenvolvimento da criança?
A Geografia \textbf{ampara} a criança para \textbf{que} ela se situe em seu espaço de convivência, a partir de sua relação com o ambiente \textbf{que} a cerca e os diversos \textbf{atores} nele envolvidos, ampliando, assim, sua noção de espaço.
Assim, a Geografia possibilitará \textbf{que} a criança compreenda a organização de sua experiência para com o espaço em \textbf{que} \textbf{vive}, a observar a natureza e a ação dos indivíduos nas formas de apropriação espacial e enxergar a paisagem.
Sua observação se estenderá a amplitudes maiores à proporção \textbf{que} sua cosmovisão se desenvolve e a visão de mundo geográfico se amplia para as relações \textbf{que} a princípio não são observáveis ou são subjetivas ao seu entendimento, como a interlocução produtiva entre cidade e campo, os movimentos de população, a produção de resíduos \textbf{sólidos}, os deslocamentos, entre outros, \textbf{que} se constituem em \textbf{uma} série de ações de pensar o espaço.
O aluno deve, portanto, reconhecer a interdependência \textbf{atual} entre campo e cidade e identificar características da produção, fluxos de matéria-prima e produtos, considerando os avanços tecnológicos, comunicacionais e informacionais, além de observar a dinâmica das indústrias, nos espaços urbanos e rurais.
A criança pode aprender no ensino de Geografia “ a entender os diferentes contextos do mundo e \textbf{que} o lugar onde ela mora não está fora do mundo, mas é parte dele ” ( SILVA \textbf{et} al., 2013:14 ).
Para tal torna -se necessário a ela observar, descrever, representar, construir explicações, criar hipóteses e compreender as diferentes manifestações da natureza.
Ao mesmo tempo, precisa saber olhar, registrar e analisar a apropriação e transformação do espaço pela ação de seu grupo social e as diferenças nos modos de outros grupos.
Para Silva \textbf{et} al. ( 2013:15 ), a Geografia só tem sentido se for pensada como instrumento para entender o espaço.
Nos anos iniciais do Ensino Fundamental, observar o cotidiano das crianças e refletir sobre essa condição ajuda a desenvolver o olhar espacial.
Indagar sobre o espaço em \textbf{que} \textbf{vivem} ( a casa, o bairro, a comunidade, a aldeia, o sítio, a chácara, a igreja, o quilombo, a vila, a cidade, o estado etc. ), significa pensar a respeito dele e, dessa forma, reunir elementos \textbf{que} possibilitam analisar suas potencialidades e demandas.
À medida \textbf{que} cresce, a criança vai desenvolvendo a capacidade de \textbf{percepção} do mundo, deixando de \textbf{depender} apenas da observação empírica e do fato concreto, para construir seus conhecimentos, sem necessariamente estar ligada a sons, cores, movimentos, formas indispensáveis para a criança pequena.
Seus conhecimentos evoluirão com suas relações sociais e a ampliação de escala do conhecimento espacial do mundo, o \textbf{que} ocorrerá em toda a sua vida.
Isso permite \textbf{que} ela leia a paisagem, o \textbf{que} deve ser explorado na sala de aula.
Nos anos iniciais, a preocupação com o entendimento dos espaços e suas articulações com os outros deve ser prática comum.
Para eles, os conceitos a serem trabalhados são aqueles \textbf{que} permitem a leitura da paisagem, fruto de \textbf{uma} dinâmica \textbf{que} reflete as coisas do lugar, sua história, seus cheiros, odores, sons, cores e movimentos.
Ler essa paisagem \textbf{implica} perceber seus elementos, oportunizar situações para \textbf{que} se \textbf{desvendem} as histórias e as relações do lugar. ( SILVA \textbf{et} al., 2013:25 ).
É necessário \textbf{que} a Geografia esteja pautada na busca dinâmica de \textbf{um} movimento vivo capaz de \textbf{aproximar} a natureza científica do cotidiano da escola e, ao mesmo tempo, construir \textbf{uma} concepção e reconstruir \textbf{uma} história, no sentido de \textbf{que} possa ser \textbf{vista} como o instrumento importante \textbf{que} é para a compreensão do mundo, traduzindo -se, nesse nível de ensino, em permitir \textbf{que} a criança desenvolva as seguintes habilidades ( PORTO, 2003:21 ) : - Compreender o seu entorno. - Descobrir como se pensa e se organiza o espaço. - Refletir sobre a relação entre a ação e o objeto. - Explicar a atuação humana. - Identificar a forma como o \textbf{homem} interage com o mundo de maneira a transformá -lo.
Para a mesma autora, a criança precisa também perceber a relação entre o \textbf{homem} e o meio ; e \textbf{que} a cada dia se torna mais importante o desenvolvimento sustentável.
Essas noções ecológicas se iniciam na infância e o aluno deve ser chamado a participar desse processo, procurando compreender \textbf{que}, na relação com o espaço, o \textbf{homem} preenche com a cultura os espaços geográficos e os tempos históricos, assumindo posturas de agente de transformação, preenchendo -o com a cultura os espaços geográficos e os tempos históricos.
O \textbf{homem} nessa trajetória se identifica com sua própria ação : objetiva o tempo, temporaliza -se, faz -se homem-histórico ( \textbf{FREIRE}, 1990:31 ).
Santos ( 1998:20 ) afirma \textbf{que} “ todos os dias o povo se \textbf{renova} ; o povo \textbf{renova} cotidianamente o seu estoque de impressões, de conhecimentos, de luta ”.
É necessário \textbf{que} a criança seja parte dessa realidade, já \textbf{que} a sociedade \textbf{atual} dá a ela \textbf{um} papel passivo, sem envolvimentos com o mundo adulto.
Para o mesmo \textbf{autor} ( 1998:51 ), a Geografia tem \textbf{um} papel importante nesse contexto, pois > à Geografia cabe estudar o conjunto indissociável de sistemas de objetos e sistemas de ação \textbf{que} formam o espaço tanto mais artificial \textbf{hoje} quanto mais estranho ao lugar.
Seus habitantes, suas fábricas e hidroelétricas são objetos \textbf{que} marcam o espaço, dando -lhe \textbf{um} conteúdo técnico.
A criança é parte dessa realidade e deve utilizar como ponto de partida as relações espaciais topológicas² locais, pois, segundo Freire ( 1990 ), a localização do \textbf{educando} é o ponto de partida para a construção do mundo. > O \textbf{que} tenho \textbf{dito} sem cansar, e redito, é \textbf{que} não podemos deixar de \textbf{lado}, desprezado como algo imprestável, o \textbf{que} educandos, sejam crianças chegando à escola ou jovens e adultos a centros de educação \textbf{popular}, trazem consigo de compreensão de mundo, das mais variadas dimensões de sua prática social.
Sua fala, sua forma de contar, de calcular, seus saberes em torno da saúde, do corpo, da \textbf{sexualidade}, da vida, da morte, da força dos \textbf{santos}, dos conjuros ( \textbf{FREIRE}, 1990:86 ) Essa compreensão do mundo, esses saberes \textbf{que} são trazidos pela criança, constituem -se como parte dessa alfabetização de leitura de mundo, objeto importante em nossas observações.
Concordamos com Perez ( 2001:107 ) \textbf{que} isso significa respeitar os saberes \textbf{populares}.
A autora explica \textbf{que} > Alfabetizar -se na leitura de mundo é respeitar esses saberes de \textbf{que} nos fala Freire, é valorizar a sabedoria \textbf{que} resulta das experiências culturais locais da criança, para possibilitar \textbf{que} ele avance para além de suas crenças em torno de si-com-o-mundo. ² Topológico \textbf{diz} respeito às posições relativas de \textbf{uma} figura geométrica.
As características \textbf{que} não são afetadas pelas alterações de tamanho e de ângulo. \textbf{Percepção} ligada ao \textbf{ver}, manipular, apalpar, vivenciar e movimentar.
As primeiras relações espaciais \textbf{que} a criança estabelece são as chamadas relações espaciais topológicas elementares, ou seja, \textbf{que} se estabelecem a partir do espaço próximo dela com referenciais como : dentro, fora, ao \textbf{lado}, na frente, atrás, perto, longe, sem considerar ainda distâncias medidas e ângulos ( ALMEIDA e PASSINI, 2002:31 ).
Para Callai ( 1998:59 ), o lugar em \textbf{que} a criança \textbf{vive} é o ponto de partida dela para a compreensão dos fenômenos e de sua explicação, pois “ por ele é mais fácil organizar as informações, podendo -se teorizar, abstrair do concreto, na busca de explicações, de comparações e de extrapolações ”.
Nesse percurso, faz -se necessário \textbf{que} os professores dos anos iniciais estejam atentos à importância da alfabetização \textbf{cartográfica}, adotando estratégias \textbf{que} possibilitem ao aluno descrever as características dos espaços mais próximos do próprio cotidiano, como os de moradia e vivência.
A identificação dos atributos e funções dos diferentes locais, em ambiente urbano e rural, público e privado, como casas, apartamentos, moradias, escolas, praças, mercados, entre outros, é indispensável para \textbf{que} ocorra o reconhecimento das semelhanças e as diferenças entre esses lugares, bem como a \textbf{percepção} das características comuns ou específicas \textbf{que} os diferenciam – por exemplo, a presença em \textbf{um} mesmo espaço de situações ou condições socioeconômicas diversas, como formas de habitações distintas, \textbf{que} \textbf{revelem} aspectos claros de desigualdades sociais, tais como : se os locais são abertos ou fechados ; se são grandes ou pequenos ; se há neles pessoas \textbf{conhecidas} ; se circulam mais crianças ou adultos, entre outros.
É relevante desenvolver o reconhecimento da importância de atitudes responsáveis com o meio onde \textbf{vive} o aluno e com o ambiente em \textbf{que} se relaciona, fazendo -o refletir sobre a necessidade de acordos para o convívio harmônico entre os grupos sociais.
A partir dessa perspectiva pode -se construir e dar significado, coletivamente, a combinados para regular os comportamentos nos diferentes espaços ( sala de aula, pátio etc. ).
Por exemplo, tratar de regras de convivência \textbf{que} sejam convergentes para comportamentos sustentáveis e ambientalmente responsáveis, como : não jogar resíduos no chão, zelar pelo bom convívio social, respeitar os colegas e os professores, entre outros.
O aluno deve conhecer e respeitar os costumes do bairro, da comunidade ou da cidade, identificando as tradições dos grupos sociais presentes no cotidiano.
Deve comparar os costumes das diferentes populações : quais as festas, feiras, manifestações culturais, comemorações \textbf{que} fazem parte da comunidade, qual a origem e/ou significado dos diversos costumes e tradições.
Ao longo do Ensino Fundamental, o aluno deve identificar e reconhecer a diversidade socioespacial existente na comunidade, comparar diferentes grupos presentes na escola e em seu entorno, assim como aspectos diversos \textbf{que} sejam pertinentes às características relacionadas à cidade, ao campo, povos indígenas, quilombolas, \textbf{ribeirinhos}, ciganos etc.
Em relação ao território maranhense, ele deve identificar e conhecer os espaços da produção agropecuária, extrativista industrial e comercial a partir do cotidiano, reconhecendo diferentes produtos e espaços de produção.
Deve, ainda, reconhecer os passos para a transformação da matéria-prima na produção de bens e serviços.
Ao término dos anos iniciais, o aluno deve ter adquirido competência para identificar e reconhecer as transformações de paisagens ocorridas no espaço e na sociedade, semelhanças, diferenças em relação a ritmos das mudanças, aspectos da estrutura, e \textbf{que} se relacionem com o cotidiano das realidades em \textbf{que} está inserido, com base em diferentes representações.
Portanto, o professor deve prosseguir coordenando a construção do conhecimento a partir do \textbf{raciocínio} geográfico \textbf{que} contemple também os diversos problemas ambientais \textbf{que} deverão ser abordados como parte de \textbf{um} processo mais amplo, complexo e interconectado ao impacto das ações humanas sobre a natureza e em todo contexto socioeconômico e cultural no qual o aluno está inserido.
Diante dos desafios a serem superados nessa etapa, o aluno deve conhecer os problemas ambientais locais, regionais e globais, tanto em relação ao campo quanto à cidade, percebendo o impacto das ações humanas sobre os componentes naturais e socioeconômicos \textbf{que} constituem esses diversos espaços.
Essa \textbf{abordagem} inicial também deve contemplar a diversidade de aspectos \textbf{que} constituem o território maranhense, desde elementos \textbf{imediatos} relacionados às áreas de entorno tanto das habitações quanto das escolas em \textbf{que} os alunos estão inseridos, como também devem evidenciar aspectos mais abrangentes \textbf{que} sejam característicos das bases físicas e antrópicas \textbf{que} integram a complexidade dos elementos e compõem o estado do Maranhão.
Como exemplo dessa condição, podemos apontar as transformações decorrentes das atividades humanas sobre a paisagem maranhense, assim como aquelas relacionadas aos processos históricos e sociais \textbf{que} elaboram o patrimônio cultural do nosso estado.
Pretende -se, portanto, \textbf{que} ao término dos anos iniciais o aluno tenha condições de refletir sobre tais especificidades e contribuir com a apresentação de propostas iniciais \textbf{que} repercutam na mitigação dos aspectos analisados, \textbf{que} a Geografia o tenha \textbf{aproximado} de sua realidade dando -lhe elementos para fazer leitura do estado do Maranhão e do seu espaço em outras escalas.
Geografia nos anos finais No Ensino Fundamental, os conhecimentos geográficos devem ser trabalhados e articulados com os saberes dos demais componentes curriculares no sentido de promover o desenvolvimento de diferentes \textbf{raciocínios} \textbf{que} contemplem os aspectos \textbf{inerentes} à Geografia ( espaço geográfico, relação temporal, aspectos socioeconômicos, culturais, etc. ).
Nos anos finais do Ensino Fundamental, o aluno se depara com grandes mudanças, sejam elas no tratamento dado aos conhecimentos a serem trabalhados nos diversos componentes curriculares, ou mesmo por questões ligadas à organização curricular, relacionadas aos tempos e espaços escolares, entre tantas outras.
O desenho da dinâmica da sala de aula passa a ser outro, com \textbf{uma} variedade maior de professores e \textbf{disciplinas}, \textbf{que} lhe \textbf{exigem} adaptação a esse novo momento.
Dessa forma, faz -se necessário \textbf{um} olhar diferenciado para a maneira como serão desenvolvidos os conhecimentos geográficos nessa fase da vida do aluno, de modo a não promover rupturas bruscas entre o tratamento dado ao conhecimento nos anos iniciais e seu desenvolvimento nos anos finais do Ensino Fundamental.
Entre o 6º e o 9º ano do Ensino Fundamental, o tratamento dado ao conhecimento em \textbf{termos} de complexidade e densidade irá sofrer alterações \textbf{que} devem acompanhar o nível de \textbf{maturidade} social e cognitiva dos alunos, de modo \textbf{que} possam realizar \textbf{uma} transição mais amena entre o Ensino Fundamental e o Ensino Médio.
Nesse contexto, a realidade local deve ser aprofundada, de modo a estimular \textbf{uma} leitura \textbf{crítico-reflexiva} das questões socioespaciais do entorno dos educandos, \textbf{que} pode acontecer a partir de múltiplas linguagens, tais como mapas, gráficos, cartogramas, infográficos, entre outros, para ampliar a leitura \textbf{cartográfica} e geográfica desde os anos iniciais.
Como exemplo, o texto O Maranhão sob perspectiva geográfica e o mapa com a visualização da organização espacial maranhense a seguir oferecem dois tipos de leituras para a construção desse conhecimento local.

% matched lemmas: abordagem, amparar, aproximar, ator, atual, autor, cartográfico, conceitual, conhecido, crítico-reflexivo, depender, desvendar, direção, disciplina, dizer, educar, et, exigir, freire, historicamente, hoje, homem, imediato, implicar, inerente, lado, maturidade, metodológico, percepção, pesquisador, popular, preciso, que, raciocínio, renovar, revelar, ribeirinho, santo, sexualidade, sujeito, sólido, termos, um, ver, viver
\end{textsample}
